
\begin{frame}{Le  fenêtrage  avec Flink}

Flink permet de découper les flux en \alert{fenêtres}. Trois types de fenêtres possibles.

\vfill

\figSlideWithSize{flink-windows}{7.5cm}

\end{frame}



\begin{frame}{Opérations sur les fenêtres}

Attention à bien évaluer les ressources nécessaires ! Il peut être nécessaire d'accumuler complètement le contenu 
d'une fenêtre avant de déclencher l'agrégation.

\vfill

\begin{redbullet}
   \item  application \alert{incrémentale} d’une fonction \textit{Reduce()} 
      combinant deux éléments de la fenêtre
   \item    application \alert{incrémentale} d’une fonction \textit{Fold()}
   ajoutant un élément de la fenêtre à un accumulateur
   \item   enfin, toute fonction s’appliquant à l’ensemble des éléments de la fenêtre, et qui ne peut pas s’appliquer incrémentalement.
\end{redbullet}

\vfill

La troisième est potentiellement consomatrice de mémoire.

\end{frame}



\begin{frame}[fragile]{Un premier exemple de fenêtre}

Au préalable, lancer notre générateur de flux sur le port 9000.

\vfill

\begin{minted}{scala}
import org.apache.flink.streaming.api.windowing.assigners._;

val stream = senv.socketTextStream("localhost", 9000, '\n')
val w = stream.map ( { x => Tuple1(x.toInt) } )
              .windowAll(TumblingProcessingTimeWindows.of(Time.seconds(5)))
               .fold("Liste: ") { (acc, v) => acc + " | " + v }
               .print()
senv.execute(" ")
\end{minted}

\vfill

On cumule les entiers de la fenêtre, toutes les 5 secondes.
\end{frame}


\begin{frame}[fragile]{Avec fenêtre glissante}

La fenêtre couvre 10 secondes, et est évaluée toutes les 5 secondes.
\vfill

\begin{minted}{scala}
import org.apache.flink.streaming.api.windowing.assigners._;

val stream = senv.socketTextStream("localhost", 9000, '\n')
val w = stream.map ( { x => Tuple1(x.toInt) } )
               .windowAll(SlidingProcessingTimeWindows.of(Time.seconds(10), Time.seconds(5)))
               .fold("Liste: ") { (acc, v) => acc + " | " + v }
               .print()
senv.execute(" ")
\end{minted}

\vfill

Le \texttt{windowAll()} indique qu'il n'y a pas de partitionnement. 
\end{frame}


\begin{frame}[fragile]{Dernier exemple: avec partitionnement}

L’exemple suivant partitionne le flux d’entiers en 2: les pairs et les impairs.

\vfill

\begin{minted}{scala}
import org.apache.flink.streaming.api.windowing.assigners._;

val stream = senv.socketTextStream("localhost", 9000, '\n')
case class MonEntier (classe: Int, valeur: Int)
val w = stream.map ( { x => x.toInt } ).map({ x => MonEntier (x % 2, x)} )
               .keyBy("classe")
               .window(TumblingProcessingTimeWindows.of(Time.seconds(5)))
               .fold("Liste: ") { (acc, v) => acc + " | " + v.valeur }
               .print()
senv.execute(" ")
\end{minted}

\vfill
\begin{blocgauche}
Le niveau de parallélisme autorisé par ce fenêtrage n’est que de deux.
\end{blocgauche}
\end{frame}


%%%%%%%
\begin{frame}[fragile]{En conclusion}

Le fenêtrage est un outil puissant pour effectuer l'analyse de flux (construction incrémentale de modèles).

\vfill

Deux obstacles à prendre en compte\,:
\begin{redbullet}
 \item Des calculs non incrémentaux sur de grandes fenêtres requièrent beaucoup de mémoire
 \item Le patitionnement de flux limite les possibilités de parallélisme.
\end{redbullet}

\vfill
\begin{blocgauche}
 À faire\,: exercices montrant le traitement de documents 
 structurés (cf. support)
 \end{blocgauche}
\end{frame}